\section{ОБЩАЯ ПОСТАНОВКА ЗАДАЧИ NMT}

\subsection{Вероятностная постановка задачи}

Машинный перевод заключается в построении текста на целевом языке по заданному тексту на исходном языке.
В нейронном машинном переводе эта задача обычно формулируется как задача условного языкового моделирования.
Перед подачей в модель текст преобразуется в последовательность токенов.
В абстрактном виде токеном будем называть элемент конечного словаря: это может быть слово, часть слова, отдельный символ или специальный служебный элемент.

Пусть $V_x$ --- словарь исходного языка, а $V_y$ --- словарь целевого языка.
Исходное предложение представим последовательностью токенов
\[
	x = (x_1, \ldots, x_m), \qquad x_i \in V_x,
\]
а перевод --- последовательностью токенов
\[
	y = (y_1, \ldots, y_n), \qquad y_t \in V_y.
\]
Задача модели с параметрами $\theta$ состоит в оценке условной вероятности $p_\theta(y \mid x)$.
Авторегрессионная декомпозиция этой вероятности имеет вид:
\[
	p_\theta(y \mid x) = \prod_{t=1}^{n} p_\theta(y_t \mid y_{<t}, x),
\]
где $y_{<t} = (y_1, \ldots, y_{t-1})$ --- уже построенная часть перевода.
На каждом шаге модель предсказывает следующий токен перевода, используя исходное предложение и предыдущие токены целевой последовательности.

Для обучения используется параллельный корпус
$D = \{(x^{(i)}, y^{(i)})\}_{i=1}^{N}$, состоящий из пар исходных предложений и эталонных переводов.
Параметры модели подбираются с помощью максимизации правдоподобия, что эквивалентно минимизации отрицательного логарифма вероятности правильного перевода:
\[
	\mathcal{L}(\theta) = - \sum_{i=1}^{N} \sum_{t=1}^{n_i}
	\log p_\theta\left(y_t^{(i)} \mid y_{<t}^{(i)}, x^{(i)}\right).
\]
Такая функция потерь задает штраф за низкую вероятность токенов, которые присутствуют в эталонном переводе.
После обучения перевод для нового предложения можно записать как задачу поиска наиболее вероятной последовательности:
\[
	\hat{y} = \arg\max_y p_\theta(y \mid x).
\]
На практике точный перебор всех вариантов невозможен, поэтому используют приближенные методы декодирования, например жадный поиск или beam search.

\subsection{Токенизация текста}

Нейронная модель не работает непосредственно со строкой текста.
Перед обучением каждое предложение преобразуется в последовательность индексов словаря.
Выбор способа токенизации влияет на длину последовательностей, размер словаря и способность модели обрабатывать редкие слова.
По этой причине в машинном переводе часто используют подсловные токенизаторы, которые разбивают редкие слова на более частые фрагменты.

Одним из наиболее распространенных подходов является BPE (Byte Pair Encoding)~\cite{bpe}.
Изначально BPE был предложен как метод сжатия данных, но в нейронном машинном переводе он используется для построения подсловного словаря.
Алгоритм начинает с разбиения слов на отдельные символы, после чего многократно объединяет наиболее частую пару соседних элементов.
После заданного числа объединений получается словарь подсловных единиц.
В результате частотные слова могут кодироваться одним токеном, а редкие слова --- последовательностью более мелких фрагментов.

Другой часто используемый токенизатор называется WordPiece~\cite{wordpiece}.
Он похож на BPE тем, что также строит словарь с помощью последовательного объединения более мелких единиц в более крупные.
Однако критерий выбора объединения связан не только с частотой пары, но и с тем, насколько добавление нового токена улучшает вероятностную модель корпуса $\mathcal{C}$.
В общем виде очередное объединение можно описать как выбор пары, дающей наибольший прирост логарифма правдоподобия:
\[
	(a, b)^* = \arg\max_{(a,b)} \Delta \log P(\mathcal{C}).
\]
На практике WordPiece часто помечает токены, продолжающие слово, специальным префиксом.
За счет этого начало слова отделяется от его внутренней части.

Токенизатор Unigram использует другой принцип, не использующий слияния подслов~\cite{unigram}.
В нем заранее задается большой набор возможных подсловных единиц, а затем обучается вероятностная модель, в которой предложение рассматривается как результат выбора последовательности подсловных токенов.
Если $u = (u_1, \ldots, u_k)$ --- один из возможных вариантов разбиения предложения на подслова, то его вероятность можно записать как
\[
	P(u) = \prod_{j=1}^{k} p(u_j).
\]
Из всех допустимых разбиений выбирается наиболее вероятное:
\[
	u^* = \arg\max_{u \in S} P(u),
\]
где $S$ --- множество возможных сегментаций исходной строки.
В процессе обучения маловероятные подсловные единицы постепенно удаляются из словаря.
Преимущество Unigram --- возможность задать несколько допустимых разбиений одного и того же текста и использовать это для регуляризации обучения.

\subsection{Метрики качества перевода}

Качество системы машинного перевода обычно оценивают сравнением с одним или несколькими эталонными переводами.
Пусть $\hat{y}$ --- перевод, построенный моделью, а $r$ --- эталонный перевод.
Автоматическая метрика должна численно оценить, насколько $\hat{y}$ близок к $r$ по форме или по смыслу.

Классической метрикой машинного перевода является BLEU~\cite{bleu}.
Она основана на точности совпадения $n$-грамм между переводом модели и эталоном.
Для каждого порядка $n$ вычисляется модифицированная точность $p_n$, в которой число совпавших $n$-грамм ограничивается их количеством в эталонном переводе.
Итоговая оценка записывается как
\[
	\operatorname{BLEU} = BP \cdot \exp\left(\sum_{n=1}^{N} w_n \log p_n\right),
\]
где $w_n$ --- веса разных порядков $n$-грамм, а $BP$ --- штраф за слишком короткий перевод.
Обычно используют $N = 4$ и равные веса.
Метрика BLEU удобна для сравнения систем на корпусном уровне, но плохо учитывает синонимы и другие допустимые переформулировки.

Метрика chrF, или CHRF, оценивает совпадение символьных $n$-грамм~\cite{chrf}.
В отличие от BLEU, она учитывает не только точность, но и полноту.
Пусть $P_{chr}$ --- средняя точность по символьным $n$-граммам, а $R_{chr}$ --- средняя полнота.
Тогда оценка имеет вид:
\[
	\operatorname{chrF}_{\beta} =
	\frac{(1 + \beta^2) P_{chr} R_{chr}}{\beta^2 P_{chr} + R_{chr}}.
\]
Так как сравнение выполняется на уровне символов, chrF лучше реагирует на частичные совпадения словоформ.
Это особенно важно для языков с развитой морфологией, где разные формы одного слова могут иметь общие символьные фрагменты.

Более современным подходом является COMET~\cite{comet}.
В отличие от BLEU и chrF, эта метрика является обучаемой нейронной моделью.
Она получает на вход исходное предложение, перевод модели и эталонный перевод, после чего предсказывает оценку качества:
\[
	q = f_{\phi}(x, \hat{y}, r),
\]
где $f_{\phi}$ --- обученная модель оценки, а $q$ --- предсказанная численная оценка.
COMET обучается на данных с человеческими оценками переводов, поэтому может учитывать смысловую близость даже в тех случаях, когда поверхностное совпадение слов невелико.
При этом такая метрика сложнее в использовании, так как требует отдельной предобученной модели и зависит от данных, на которых эта модель была обучена.

\subsection{Методы генерации перевода}

После обучения авторегрессионная модель задает распределение вероятностей следующего токена $p_\theta(y_t \mid y_{<t}, x)$. Для получения конкретного перевода необходимо выбрать правило декодирования. Детерминированные методы, такие как жадный поиск или beam search, выбирают наиболее вероятные продолжения и обычно применяются при итоговой оценке качества.
Для reinforcement learning техник (подробнее в следующем разделе) требуется получать несколько различных вариантов ответа, поэтому чаще используется стохастическая генерация.
Несмотря на то, что стохастическая генерация текста более применима к LLM моделям, которым нужно быть ``творческими'', такой метод применим и к NMT задачам, где ``разброс'' оптимального ответа куда ниже.

Одним из параметров стохастической генерации является temperature. Пусть $z_i$ --- logit токена $i$ перед softmax. При температуре $T$ распределение следующего токена имеет вид:
\[
	p_i(T) = \frac{\exp(z_i / T)}{\sum_j \exp(z_j / T)}.
\]
При $T < 1$ распределение становится более резким, и модель чаще выбирает наиболее вероятные токены. При $T > 1$ распределение сглаживается, поэтому генерация становится более разнообразной, но возрастает риск ошибок.

Другой распространенный прием --- nucleus sampling, или top-p sampling~\cite{nucleus_sampling}. На каждом шаге выбирается минимальное множество наиболее вероятных токенов $S_p$, сумма вероятностей которых не меньше заданного порога $p$:
\[
	S_p = \min \left\{ S: \sum_{i \in S} p_i \geq p \right\}.
\]
После этого вероятности токенов вне $S_p$ зануляются, а оставшееся распределение нормируется заново. В отличие от top-k sampling, где число допустимых токенов фиксировано, top-p динамически меняет размер множества кандидатов в зависимости от уверенности модели.

Как правило, temperature и top-p sampling используют вместе, чтобы ограничить маловероятные токены и при этом сохранить разнообразие генерации. Фрагмент реализации такого выбора следующего токена приведен в приложении~А.

\subsection{Методы reinforcement learning для NMT}

После обучения по параллельному корпусу модель перевода можно дополнительно оптимизировать не только по token-level cross entropy, но и по внешней оценке качества построенного перевода. В такой постановке генерация перевода рассматривается как последовательное принятие решений: на каждом шаге модель выбирает следующий токен, а после получения полной последовательности получает reward за качество результата.

Пусть $\hat{y}$ --- перевод, сгенерированный моделью для исходного предложения $x$, а $r$ --- эталонный перевод. Reward-функцию можно задать через автоматические метрики качества:
\[
	R(x, \hat{y}, r) = \alpha_1 \operatorname{BLEU}(\hat{y}, r) +
	\alpha_2 \operatorname{chrF}(\hat{y}, r) +
	\alpha_3 \operatorname{COMET}(x, \hat{y}, r),
\]
где $\alpha_1$, $\alpha_2$ и $\alpha_3$ задают вклад отдельных метрик. В простейшем случае в качестве reward можно использовать одну метрику, например COMET, либо взвешенную комбинацию BLEU, chrF и COMET.

Одним из методов такой оптимизации является GRPO (Group Relative Policy Optimization)~\cite{grpo}. Для одного входного предложения модель генерирует группу вариантов перевода, обычно с помощью стохастического декодирования с temperature и top-p, после чего каждый вариант оценивается reward-функцией. Разные варианты внутри группы дают локальное сравнение качества переводов для одного и того же исходного предложения. Вместо обучения отдельной value-модели GRPO сравнивает reward отдельного ответа со средним reward внутри группы. Если вариант лучше среднего, его вероятность увеличивается; если хуже среднего, уменьшается. Такая интерпретация делает метод удобным для задач, где качество ответа можно оценить внешней метрикой, но трудно задать правильное действие на каждом отдельном шаге генерации.

В упрощенном виде оптимизируемую функцию можно записать как clipped policy-gradient loss с KL-регуляризацией относительно опорной модели:
\[
	\mathcal{L}_{GRPO}(\theta) = -\mathbb{E}_{x, \hat{y}}
	\left[
		\min\left(\rho_\theta A, \bar{\rho}_\theta A\right)
		- \beta D_{KL}(\pi_\theta \| \pi_{ref})
		\right].
\]
Здесь $\rho_\theta = \pi_\theta(\hat{y} \mid x) / \pi_{old}(\hat{y} \mid x)$ --- отношение вероятностей ответа в новой и старой политиках, $\bar{\rho}_\theta = \operatorname{clip}(\rho_\theta, 1 - \varepsilon, 1 + \varepsilon)$ --- ограниченное отношение вероятностей, $A$ --- групповая advantage-оценка, вычисленная из reward вариантов внутри одной группы, $\varepsilon$ --- параметр ограничения шага обновления, $\beta$ --- вес KL-штрафа, а $\pi_{ref}$ --- опорная модель. Знак минус связан с тем, что при обучении обычно минимизируется loss, хотя сама policy-gradient часть соответствует максимизации ожидаемого reward. Фрагмент реализации GRPO loss приведен в приложении~А.

BLEU можно использовать как часть reward, но он плохо подходит в качестве единственного сигнала оптимизации. Метрика основана на поверхностном совпадении $n$-грамм, слабо учитывает допустимые переформулировки и на уровне отдельных предложений дает шумный и разреженный сигнал. В исследованиях reinforcement learning для NMT также отмечалось, что оптимизация непосредственно под BLEU может давать нестабильное обучение и не всегда приводит к улучшению человеческой оценки перевода~\cite{rl_nmt_bleu_reward}. Поэтому для reward в переводе предпочтительно использовать более устойчивые сигналы, например chrF и COMET, либо их комбинацию с BLEU.

\subsection{Методы регуляризации моделей NMT}

При обучении моделей NMT число параметров обычно велико, а параллельный корпус ограничен по размеру и качеству.
Поэтому модель может не только изучать общие закономерности перевода, но и запоминать частные особенности обучающих примеров.
Для уменьшения такого эффекта используют методы регуляризации: они изменяют процесс обучения так, чтобы модель давала более устойчивые предсказания на новых данных.

Одним из распространенных методов является dropout~\cite{dropout}.
Пусть \mbox{$h$~---~вектор активаций} некоторого слоя, $h = (h_1, \ldots, h_d)$, а \mbox{$p$~---~вероятность} зануления элемента.
Во время обучения случайно выбирается маска
\[
	m_i \sim \operatorname{Bernoulli}(1 - p), \qquad i = 1, \ldots, d.
\]
После этого вместо исходного вектора $h$ в следующий слой передается вектор
\[
	\tilde{h} = \frac{m \odot h}{1 - p},
\]
где $\odot$ обозначает покомпонентное умножение.
Такое масштабирование сохраняет математическое ожидание активаций: $\mathbb{E}\tilde{h}_i = h_i$.
На каждой итерации обучения сеть фактически работает с разными подмоделями, из-за чего отдельные нейроны меньше зависят от фиксированного набора соседних признаков.
В NMT dropout применяют внутри слоев модели: к embeddings, выходам attention и feed-forward блоков, а также к промежуточным представлениям между слоями.

Другой подход связан с ограничением величины весов модели.
Если параметры модели обозначить через $w$, то при обычной максимизации правдоподобия выбираются веса, при которых велика вероятность обучающих данных $p(D \mid w)$.
В байесовской постановке рассматривается другая величина --- апостериорная плотность весов при условии данных:
\[
	p(w \mid D) = \frac{p(D \mid w)p(w)}{p(D)}.
\]
Так как $p(D)$ не зависит от $w$, получается оценка максимума апостериорной вероятности (maximum a posteriori, MAP)
\[
	\hat{w}_{MAP} = \arg\max_w \left(\log p(D \mid w) + \log p(w)\right).
\]
В MAP-подходе максимизируется апостериорная плотность весов $p(w \mid D)$, а не только правдоподобие данных $p(D \mid w)$.
Если в качестве априорного распределения взять нормальное распределение $p(w) = \mathcal{N}(0, \sigma^2 I)$, то
\[
	\log p(w) = -\frac{1}{2\sigma^2}\|w\|_2^2 + C.
\]
Отсюда получается функция потерь с $\ell_2$-регуляризацией~\cite{dlbook}:
\[
	\mathcal{L}_{reg}(w) = \mathcal{L}(w) + \frac{\lambda}{2}\|w\|_2^2.
\]
В этом случае большие веса штрафуются квадратичной добавкой и заранее считаются менее предпочтительными.
В градиентном спуске это приводит к обновлению
\[
	w_{t+1} = (1 - \eta\lambda)w_t - \eta \nabla_w \mathcal{L}(w_t),
\]
где $\eta$ --- шаг обучения.
Первый множитель постепенно уменьшает веса, поэтому в оптимизаторах этот прием называют weight decay.
В современных адаптивных оптимизаторах weight decay часто реализуют как отдельный шаг уменьшения весов, а не как часть градиента функции потерь~\cite{adamw}.

Еще одним способом регуляризации является label smoothing~\cite{label_smoothing}.
В стандартном обучении правильный токен кодируется one-hot распределением: вся вероятность относится к одному элементу словаря.
Для NMT такая постановка бывает слишком жесткой, поскольку один и тот же фрагмент текста может иметь несколько допустимых переводов.
Пусть $K = |V_y|$ --- размер целевого словаря, а $y_t$ --- правильный токен на шаге $t$.
При label smoothing целевое распределение заменяется сглаженным распределением:
\[
	\tilde{q}_t(k) = (1 - \varepsilon)\mathbf{1}\{k = y_t\} + \frac{\varepsilon}{K},
\]
где $\varepsilon$ --- коэффициент сглаживания.
Функция потерь тогда записывается как
\[
	\mathcal{L}_{LS} = -\sum_t \sum_{k=1}^{K} \tilde{q}_t(k)
	\log p_\theta(k \mid y_{<t}, x).
\]
Label smoothing ограничивает чрезмерную уверенность модели в единственном токене и обычно делает распределение выходных вероятностей более устойчивым.
